\documentclass[12pt]{article}
\usepackage[utf8]{inputenc}
\usepackage[T1]{fontenc}
\usepackage{hyperref}
\title{Nyiyaparli Language Resources for K-12 Curriculum Development}
\author{Compiled by Claude Sonnet 5.5 (claude-sonnet-5-5) for Andr\'as Kornai}
\date{}
\begin{document}
\maketitle

\section{Introduction}
Nyiyaparli (ISO 639-3 xny; Pama-Nyungan; Pilbara region, Western Australia) is a severely endangered Aboriginal language with very few remaining first-language speakers. No school programme teaching it was found in the searches. The materials below are what turned up in web search. The Wangka Maya Pilbara Aboriginal Language Centre is the main institutional source. Availability and access conditions were not checked, and cultural protocols may apply, so the community and the language centre should be consulted before any use.

\section{Primary Educational Resources}
No primer, textbook or school curriculum for Nyiyaparli was found in search results. The Wangka Maya page reports a book of short stories by Nyiyaparli people living in Marble Bar. It could be a starting point for reading material, but its availability was not verified.
\begin{itemize}
\item Wangka Maya Pilbara Aboriginal Language Centre, Nyiyaparli language page: \url{https://www.wangkamaya.org.au/pilbara-languages/nyiyaparli}
\item Wangka Maya, Nyiyaparli cultural introduction: \url{https://www.wangkamaya.org.au/pilbara-languages/nyiyaparli/01-cultural-introduction}
\end{itemize}

\section{Linguistic and Reference Materials}
\begin{itemize}
\item Wangka Maya Nyiyaparli dictionary. Editions are reported for 2004 (interim dictionary and finderlists), 2008 and 2012 (adds an English--Nyiyaparli wordlist and topic wordlists). The 2012 edition treats Palykutharri, Nankilakuthu and Martuyitha as dialects. Current availability was not verified. Description at: \url{https://www.wangkamaya.org.au/pilbara-languages/nyiyaparli}
\item Jacqueline Battin's thesis on Nyiyaparli grammar, based on narratives recorded by Carl von Brandenstein in the 1960s. The search results showed an ANU Open Research repository record, which may be the related item (not confirmed): \url{https://openresearch-repository.anu.edu.au/items/e4a2bcc6-cda7-4380-af64-174b486fa7d4}
\item AIATSIS Austlang entry: \url{https://collection.aiatsis.gov.au/austlang/language/a50}
\item Wikipedia, ``Niabali'' (Nyiyaparli): \url{https://en.wikipedia.org/wiki/Niabali}
\item Native title claim description for Nyiyaparli, with historical and community background: \url{https://www.drbilldayanthropologist.com/resources/NYIYAPARLI CLAIM DESCRIPTION 030903.pdf}
\end{itemize}

\section{Religious and Community Resources}
No scripture translations were found. The Wangka Maya Pilbara Aboriginal Language Centre in South Hedland is the community language organisation to approach: \url{https://www.wangkamaya.org.au/pilbara-languages/nyiyaparli}

\section{Digital and Audio Resources}
The Wangka Maya page mentions recordings made since 1969. Details, access and permissions were not verified: \url{https://www.wangkamaya.org.au/pilbara-languages/nyiyaparli}

\section{Curriculum Development Context}
Very little published material exists, and there are only a handful of speakers. A realistic curriculum would be a heritage and revitalisation programme built with Wangka Maya and Nyiyaparli families. It would draw on the dictionary, the archival recordings and the Marble Bar stories, and would handle those materials according to community protocols. English-medium teaching of Nyiyaparli words and culture is more feasible than immersion.
\end{document}